\section{Related Work}
\label{sec:related-work}

\paragraph{Behavioral evaluation and counterfactual tests.}
Behavioral testing organizes evaluation around capabilities and expected responses to controlled changes. CheckList combines linguistic capabilities with minimum functionality, invariance, and directional tests \citep{ribeiro2020checklist}. Contrast sets probe local decision boundaries through small, meaningful edits that typically change the reference label \citep{gardner2020contrastsets}. Dynabench develops this diagnostic perspective through human-and-model-in-the-loop data collection \citep{kiela2021dynabench}. These approaches motivate evaluating the relationship between an input change and its required response. Our five-category taxonomy applies that perspective to medical tasks: it records the changed condition, the requested decision, and the evidence needed to justify revision or preservation. Categories can overlap when one case requires several operations. This gives a shared interpretation to tests drawn from sources with different task formats.

\paragraph{Medical counterfactual and robustness benchmarks.}
Medical datasets expose several distinct consequences of perturbing an input. MedEinst changes discriminative evidence to test differential diagnosis \citep{chen2026medeinst}, while MedPIC-Bench provides patient-specific medication questions with counterfactual task families \citep{medpicbench}. HIV treatment scenarios examine causal consistency across interventions, interactions, and side effects \citep{prosperi2026hiv}. Counterfactual Patient Variations and DiversityMedQA examine demographic sensitivity in clinical decisions \citep{benkirane2025bias,rawat2024diversitymedqa}; cultural-cue perturbations and AMQA extend this line to cultural context and adversarial medical questions \citep{rezaei2026cultural,xiao2025amqa}. MedPerturb varies gender cues, language style, and presentation format \citep{gourabathina2025medperturb}. MedCounterFact studies how models follow supplied medical evidence when its content conflicts with plausibility or safety \citep{mo2026faithfulness}. MedRGB and BioRAB assess retrieval conditions, including misleading or counterfactual evidence \citep{ngo2024medrgb,li2025biorab}. We bring eligible instances from these eleven sources into one benchmark while preserving their native questions, answer spaces, and reference relations. The resulting functional labels connect heterogeneous perturbations to the reasoning operation that each instance actually tests.

\paragraph{Medical RAG.}
Retrieval-augmented generation combines a language model with external information \citep{lewis2020rag}. MIRAGE and the MedRAG toolkit examine how corpora, retrievers, and language models affect medical question answering \citep{xiong2024medrag}. MedRGAG connects retrieved and parametric knowledge \citep{li2026medrgag}; MA-RAG uses multiple rounds of agentic retrieval and reasoning \citep{wu2026marag}. MedGENIE supplies generated contexts for medical questions \citep{frisoni2024medgenie}, and TC-RAG maintains a memory stack to control iterative retrieval \citep{jiang2025tcrag}. These methods provide complementary ways to acquire and organize knowledge. Our seven evaluated configurations use them to diagnose which functional requirements remain difficult, then test operation-specific heads on the resulting outputs and admissible evidence. This connects system-level evaluation to targeted repairs such as recalculating support and comparing diagnostic candidates.

\paragraph{Agent aggregation.}
Self-consistency aggregates answers across sampled reasoning paths \citep{wang2023selfconsistency}, while Mixture-of-Agents passes multiple agents' outputs between successive layers \citep{wang2025moa}. Contextual calibration addresses prompt-dependent answer preferences \citep{zhao2021calibrate}. These ideas motivate distinct controls for additional sampling, candidate comparison, and aggregation. MedGuideX transforms executable guideline logic into factual and counterfactual supervision for post-training medical models \citep{shen2026medguidex}. Our support agents execute guideline logic during inference to recompute a current answer. CF-MoA gives its five agents explicit functional responsibilities and combines their outputs through a shared input contract and routing interface. Our evaluation follows both the gains from integrating these specialized operations and the incremental effects of candidate verification, tying each system change to the behavior it is intended to improve.
